\section{Arithmetic decomposition and exact Poisson density}
\label{sec:poisson}

After inserting \Cref{lem:afe} into \eqref{eq:zeta-star}, the summands in the
direct AFE piece, before the outer coefficient weight, are
\begin{equation}\label{eq:direct-summand}
 l^{-1/2-\alpha}m^{-1/2-\beta}n^{-1/2-\gamma}
 \left(\frac{hn}{klm}\right)^{it}
 \V^+(l/Y,mn,t).
\end{equation}
The dual piece is obtained by
\begin{equation}\label{eq:dual-substitution}
 (\beta,\gamma,\V^+)\mapsto(-\gamma,-\beta,\V^-)
 \quad\text{and multiplication by }X_{\beta,\gamma}(t).
\end{equation}
The recovery estimate \eqref{eq:zeta-star-recovery} permits us to work with
these sums and return to the original cubic moment with an arbitrary power
saving.

Put
\begin{equation}\label{eq:r-equation}
                         hn-klm=r.
\end{equation}
The terms with $r=0$ are the direct and dual diagonals.  Suppose $r\neq0$ and
set
\begin{equation}\label{eq:gcd-decomposition}
 d=(h,km),\qquad h=dq,\qquad km=dv,\qquad r=d\rho.
\end{equation}
Solutions exist only if $d\mid r$.  Since $(q,v)=1$, the equation becomes
\begin{equation}\label{eq:l-congruence}
                       l\equiv-\rho\overline v\pmod q.
\end{equation}
Thus every inverse is taken modulo a coprime integer, and Poisson summation in
the fixed $l$-orientation has exact density
\begin{equation}\label{eq:poisson-density}
                             \frac1q=\frac dh.
\end{equation}
We use this same orientation for both zero and nonzero frequencies.

More explicitly, if $F$ is a smooth function of $l$, then
\begin{equation}\label{eq:poisson-formula}
 \sum_{l\equiv-\rho\overline v\ (q)}F(l)
 =\frac1q\sum_{a\in\mathbb Z}
 e\!\left(-\frac{a\rho\overline v}{q}\right)
 \int_{\R}F(y)e\!\left(-\frac{ay}{q}\right)\dd y,
 \qquad e(x)=e^{2\pi ix}.
\end{equation}
After the change of variables $x=kmy$, the combined Jacobian and progression
density is
\begin{equation}\label{eq:poisson-jacobian}
                           \frac{1}{qkm}=\frac{d}{hkm}.
\end{equation}

For later use we record the exact zero mode, including the discrete endpoints.
Choose $U\in C^\infty(\R)$ with $U(x)=0$ for $x\leq1/2$ and $U(x)=1$ for
$x\geq1$.  Multiplication by $U(l)U(n)$ does not change the discrete sum.
With
\begin{equation}\label{eq:xy-def}
                         y=kml,\qquad x=hn=y+r,
\end{equation}
the direct zero mode for fixed $h,k$ is
\begin{align}
 \mathcal O^+_{0,Y}(h,k)
 ={}&\int_{\R}w(t)\sum_{m\geq1}
 \sum_{\substack{r\neq0\\d\mid r}}
 d\,h^{-1/2+\gamma}k^{-1/2+\alpha}m^{-1+\alpha-\beta}
 \notag\\
 &\quad\times\int_{\max(0,-r)}^\infty
 y^{-1/2-\alpha-it}x^{-1/2-\gamma+it}
 \V^+\!\left(\frac{y}{kmY},\frac{mx}{h},t\right)\notag\\
 &\qquad\times
 U\!\left(\frac{y}{km}\right)U\!\left(\frac{x}{h}\right)
 \dd y\dd t.
 \label{eq:exact-zero-mode}
\end{align}
The outside factor $a_h\overline{b_k}/\sqrt{hk}$ is not included in
\eqref{eq:exact-zero-mode}.  Formula \eqref{eq:exact-zero-mode} follows
directly from \eqref{eq:poisson-density} after changing $\dd l$ to
$\dd y/(km)$; in particular, the factor $d$ has not been suppressed.  The
dual zero mode is obtained by \eqref{eq:dual-substitution}.

Let $\widetilde{\mathcal O}_{0,Y}^+$ be
\eqref{eq:exact-zero-mode} with both $U$ factors deleted, initially wherever
the resulting integrals and sums converge absolutely and subsequently by the
continuation constructed in \Cref{sec:endpoint}.  Define
$\widetilde{\mathcal O}_{0,Y}^-$ in the same way.

For the nonzero-frequency analysis, decompose all five arithmetic variables
into dyadic blocks
\[
 h\asymp H_0,\quad k\asymp K_0,\quad l\asymp L_0,\quad
 m\asymp M_0,\quad n\asymp N_0
\]
and put
\begin{equation}\label{eq:S-def}
                         S=H_0N_0+K_0L_0M_0.
\end{equation}
Repeated integration by parts in $t$ shows that a non-negligible block must
satisfy
\begin{equation}\label{eq:r-localization}
 H_0N_0\asymp K_0L_0M_0,\qquad
 |r|\ll ST^{-1+\eps}.
\end{equation}
Consequently
\begin{equation}\label{eq:S-balance}
 S\asymp H_0N_0\asymp K_0L_0M_0
  \asymp(H_0K_0L_0M_0N_0)^{1/2}.
\end{equation}
The AFE conductor also gives
\begin{equation}\label{eq:MN-conductor}
                         M_0N_0\ll T^{1+\eps},
\end{equation}
while the long smoothing has $L_0\ll YT^\eps$.  The logarithmic number of
blocks is absorbed into $T^\eps$.
